\section{The experiment and the spectral metric}\label{sec:model}
Fix $D>0$, a degree $d\ge1$, and $0<\pi<1/2$. A closed binary model
point consists of two monic real polynomials $f_1,f_2$ of degree $d$,
all of whose roots lie in $[0,D]$, weights $\alpha_1+\alpha_2=1$ with
$\alpha_a\ge\pi$, and column-stochastic matrices
$U=(u_1,u_2)$ and $V=(v_1,v_2)$ with nonnegative entries. Set
\begin{equation}\label{eq:model}
 K(z)=\sum_{a=1}^2\alpha_af_a(z)u_av_a^{\mathsf T},\qquad
 q(z)=\one^{\mathsf T}K(z)\one=\sum_{a=1}^2\alpha_af_a(z),\qquad
 P_j=\frac{K(T_j)}{q(T_j)}.
\end{equation}
The clocks satisfy $D<T_1<\cdots<T_\ell$, and $w_j>0$, $\sum_jw_j=1$.
Consequently $q(T_j)>0$, including at every closed-model competitor.
The target polynomial is $A=f_1f_2$. The model is compact in monic
coefficients and stochastic entries. Component permutations are identified.
The matrices, roots, clocks and design weights are allowed to vary when a
family of centres is considered, but every competitor uses its centre's
fixed design.

For probability arrays define
\[
 h_w(P,Q)^2=\sum_jw_j\sum_e(\sqrt{P_{je}}-\sqrt{Q_{je}})^2.
\]
For two multisets $Z=\{z_i\}_{i=1}^n$ and $Z'=\{z'_i\}_{i=1}^n$, put
\[
 d_\infty(Z,Z')=\min_{\sigma\in\mathfrak S_n}\max_i|z_i-z'_{\sigma(i)}|.
\]
For real roots this is the maximum discrepancy of their increasing
orderings. All diameters in this article are diameters of sets of
multisets, not maximal distances from the centre. Write
\begin{equation}\label{eq:modulus}
 \omega_\theta(t)=\diam\{Z(A_\xi):h_w(F_T(\xi),P_\theta)\le t\}.
\end{equation}
A local modulus $\omega_{\theta,\loc}$ restricts the competitors to a
specified coefficient neighbourhood of the centre's component pair.
This restriction is never implicit.

At a strictly positive datum the Fisher inner product on zero-sum arrays is
\[
 \langle H,H'\rangle_{I,\theta}=\sum_{j,e}\frac{w_j}{P_{\theta,je}}H_{je}H'_{je}.
\]
For component increments $p_a$ of degree below $d$ and stochastic
increments $\eta=(\dot\alpha_1,\dot u_1,\dot u_2,\dot v_1,\dot v_2)$,
where channel increments have entry sum zero, define
\begin{align}
 \dot K&=\sum_a\bigl\{\alpha_ap_au_av_a^{\mathsf T}
 +f_a(\dot\alpha_au_av_a^{\mathsf T}
       +\alpha_a\dot u_av_a^{\mathsf T}+\alpha_au_a\dot v_a^{\mathsf T})\bigr\},\nonumber\\
 \dot q&=\sum_a(\alpha_ap_a+\dot\alpha_af_a),\qquad
 S_\theta(p_1,p_2,\eta)_j
       =\frac{\dot K(T_j)-P_{\theta,j}\dot q(T_j)}{q(T_j)}.
                                                        \label{eq:score}
\end{align}
Here $\dot\alpha_2=-\dot\alpha_1$. These rational formulas define every
Fisher program below in the original experiment coordinates.

\begin{lemma}[Normalizer recovery]\label{lem:normalizer}
Suppose $f_1,f_2$ are coprime, $V$ is invertible, the weights are positive,
and $\ell\ge2d+1$. The observed array determines $q$ and $K$ uniquely
among all closed-model competitors. Their coefficients are locally
Lipschitz functions of the observations, with uniform constants on
compact sets of such centres. No rank assumption on $U$ is required.
\end{lemma}
\begin{proof}
Let $e$ have degree below $d$, and let $H$ be a matrix polynomial of
degree at most $d$ satisfying $H(T_j)=e(T_j)P_j$. Interpolation of the
polynomial $qH-eK$, of degree at most $2d$, gives $qH=eK$.
Its second marginal, multiplied by
$\diag(\alpha_1,\alpha_2)^{-1}V^{-1}$, gives $q\mid ef_1$ and
$q\mid ef_2$. Coprimality implies $q\mid e$, hence $e=0$ and $H=0$.

More explicitly, let $\Pi_T$ project onto the orthogonal complement of
the degree-$d$ Vandermonde column space. The map
\[
 N_\theta:e\longmapsto\Pi_T(e(T_j)P_{\theta,j})_j,\quad\deg e<d,
\]
is injective. For a competitor, subtraction of the normalization
equations yields
$\|N_\theta(q'-q)\|\le C\max_j\|P'_j-P_j\|$; the competitor
normalizer coefficients are bounded by compactness of the root interval.
The smallest singular value of $N_\theta$ is positive. A bounded left
inverse recovers $q'-q$, and Vandermonde interpolation recovers $K'-K$.
The singular values and the coefficient bounds have positive uniform
margins on compact sets of the stated centres.
\end{proof}

\begin{lemma}[Full-rank coefficient inverse]\label{lem:fullrank}
Under Lemma~\ref{lem:normalizer}, suppose also that $U$ is invertible.
Then every sufficiently close observation has, up to one component
permutation, coefficient and stochastic parameter error bounded by a
constant times its observation error. Every exact competitor has the
same component polynomials and stochastic parameters. The constants
are uniform on compact subsets of these centres.
\end{lemma}
\begin{proof}
Choose two clocks such that
$f_1(T_j)/f_1(T_i)\ne f_2(T_j)/f_2(T_i)$. Otherwise the nonzero polynomial
$f_2(T_i)f_1-f_1(T_i)f_2$, of degree at most $d$, would vanish at all
$2d+1$ clocks. The matrix $K(T_j)K(T_i)^{-1}$ has those two distinct
eigenvalues and eigenvectors $u_1,u_2$. Normalize the eigenvectors to
have entry sum one. Simple-eigenvalue inversion is analytic in a
neighbourhood of the centre. Applying the recovered $U^{-1}$ to $K(z)$
and taking row sums recovers $\alpha_af_a$; leading coefficients and
normalized rows recover $\alpha_a,V$. This construction applies to
arbitrary closed-model competitors because their $K'(T_i)$ stays
invertible. It also proves exact uniqueness. A finite neighbourhood
cover gives compact-uniform constants. Equivalently, the derivative in
monic coefficient and stochastic coordinates has a bounded left inverse.
\end{proof}
